\documentclass[a4paper,12pt]{article}
\usepackage[UTF8]{ctex}
\usepackage{geometry}
\usepackage{enumitem}
\usepackage{amsmath}
\usepackage{xcolor}
\usepackage{listings}
\usepackage{tcolorbox}

\geometry{left=2.5cm, right=2.5cm, top=2.5cm, bottom=2.5cm}

\newtcolorbox{bluebox}[1][]{
colback=blue!5!white,
colframe=blue!75!black,
fonttitle=\bfseries,
title=#1,
arc=0mm,
boxrule=1pt,
left=2mm,right=2mm,top=1mm,bottom=1mm
}

\lstset{
language=C,
basicstyle=\ttfamily\small,
keywordstyle=\color{blue},
commentstyle=\color{green!60!black},
stringstyle=\color{red},
numbers=left,
numberstyle=\tiny\color{gray},
frame=single,
breaklines=true,
escapeinside={(*@}{@*)}
}

\title{\textbf{第五章 指针-题库深度解析}}
\author{}
\date{}

\begin{document}

\maketitle

\tableofcontents
\newpage

% ======== 题目1: 指针与数组 ========
\section*{题目1: 指针访问数组元素}

\textbf{原题目:}
若有定义:\texttt{int a[10], *p=a;},则 \texttt{*(p+5)} 表示\_\_\_\_\_\_。
\begin{enumerate}
	\item[A)] a[5]的地址
	\item[B)] a[5]的值
	\item[C)] a[6]的值
	\item[D)] a的地址
\end{enumerate}

\begin{bluebox}{答案与深度解析}
	\textbf{答案: B) a[5]的值}

	% ======== 阶段一:核心认知框架 ========
	\textbf{【核心认知框架】}
	\begin{itemize}[leftmargin=*, nosep]
		\item \textbf{题目本质}: 考查指针与数组的关系,区分地址与值
		\item \textbf{关键区分点}: \texttt{p+i} VS \texttt{*(p+i)} VS \texttt{\&a[i]}
		\item \textbf{命题人意图}: 测试考生对指针运算和数组访问的理解
	\end{itemize}

	% ======== 阶段二:深度解析 ========
	\textbf{【指针与数组深度解析】}

	\textbf{① 指针初始化}: \texttt{int a[10], *p=a;}
	\begin{itemize}[leftmargin=*, nosep]
		\item \texttt{a[10]}: 定义含10个int元素的数组
		\item \texttt{*p=a}: p指向数组a的首元素a[0]
		\item 等价于: \texttt{*p=\&a[0];}
	\end{itemize}

	\textbf{② 指针运算规则}
	\begin{itemize}[leftmargin=*, nosep]
		\item \texttt{p+1}: p向后移动1个int(4字节)
		\item \texttt{p+5}: p向后移动5个int(20字节)
		\item \texttt{p+i}的结果类型仍是\texttt{int*},指向\texttt{a[i]}
	\end{itemize}

	\textbf{③ 解引用(\*)的作用}
	\begin{tabular}{|c|c|c|}
		\hline
		\textbf{表达式} & \textbf{类型} & \textbf{含义}    \\
		\hline
		p            & int*        & 指针变量,指向a[0]    \\
		\hline
		p+5          & int*        & 指针,指向a[5]      \\
		\hline
		*(p+5)       & int         & a[5]的值         \\
		\hline
		\&a[5]       & int*        & a[5]的地址        \\
		\hline
		a+5          & int*        & a[5]的地址,等价于p+5 \\
		\hline
		*(a+5)       & int         & a[5]的值,等价于a[5] \\
		\hline
	\end{tabular}

	\textbf{④ 汇编级理解}
	\begin{lstlisting}[language=C]
// C代码: *(p+5)

// x86汇编:
mov eax, DWORD PTR [ebp-12]  ; 加载p的值(数组首地址)
add eax, 20                   ; p+5,向后移动5个int(5*4=20字节)
mov eax, DWORD PTR [eax]      ; 解引用,读取该位置的值

// 对比: &a[5]
lea eax, [ebx+20]             ; 计算a[5]的地址,加20字节offset

// 对比: a[5]
mov eax, DWORD PTR [ebx+20]   ; 直接读取a[5]的值
    \end{lstlisting}

	\textbf{⑤ 指针与数组的等价关系}
	\begin{lstlisting}[language=C]
// 以下写法完全等价:
int a[10], *p = a;

// 方法1: 下标法
a[5]    // 值
&a[5]   // 地址

// 方法2: 指针法
*(p+5)  // 值
p+5     // 地址

// 方法3: 数组名法
*(a+5)  // 值
a+5     // 地址
    \end{lstlisting}

	% ======== 选项分析 ========
	\textbf{【选项分析】}

	\textbf{A) a[5]的地址 — 错误}
	\begin{itemize}[leftmargin=*, nosep]
		\item 这是\texttt{p+5}或\texttt{\&a[5]}的含义
		\item \texttt{*(p+5)}是地址的内容(值),不是地址本身
	\end{itemize}

	\textbf{B) a[5]的值 — 正确}
	\begin{itemize}[leftmargin=*, nosep]
		\item \texttt{p+5}得到a[5]的地址(类型:int*)
		\item \texttt{*(p+5)}对该地址解引用,得到a[5]的值(类型:int)
		\item 这与\texttt{a[5]}完全等价
	\end{itemize}

	\textbf{C) a[6]的值 — 错误}
	\begin{itemize}[leftmargin=*, nosep]
		\item a[6]的值应该是\texttt{*(p+6)}
		\item \texttt{*(p+5)}是a[5],不是a[6]
	\end{itemize}

	\textbf{D) a的地址 — 错误}
	\begin{itemize}[leftmargin=*, nosep]
		\item a的地址是\texttt{p}或\texttt{\&a[0]}
		\item 这里问的是\texttt{*(p+5)},明显不是数组首地址
	\end{itemize}

	% ======== 阶段三:应背诵知识点 ========
	\textbf{【应背诵知识点】}

	\begin{enumerate}[leftmargin=*, nosep]
		\item \textbf{指针运算公式}:
		      \begin{itemize}
			      \item \texttt{p+i} = \texttt{\&a[i]}
			      \item \texttt{*(p+i)} = \texttt{a[i]}
			      \item \texttt{\&*(p+i)} = \texttt{p+i} = \texttt{\&a[i]}
		      \end{itemize}

		\item \textbf{三种访问方式}:
		      \begin{itemize}
			      \item \textbf{下标法}: \texttt{a[i]}, \texttt{p[i]}
			      \item \textbf{指针法}: \texttt{*(a+i)}, \texttt{*(p+i)}
			      \item \textbf{混合法}: \texttt{a[i]}, \texttt{*(p+i)}
		      \end{itemize}

		\item \textbf{地址与值}:
		      \begin{itemize}
			      \item \texttt{p+i}是地址
			      \item \texttt{*(p+i)}是值
			      \item \texttt{\&*(p+i)} = 地址
		      \end{itemize}
	\end{enumerate}

	% ======== 阶段四:实战策略 ========
	\textbf{【实战策略】}
	\begin{itemize}[leftmargin=*, nosep]
		\item \textbf{分析步骤}:
		      \begin{enumerate}
			      \item 确定指针p指向哪里
			      \item p+5指向哪个元素
			      \item *的作用是取值还是取地址
		      \end{enumerate}

		\item \textbf{记忆口诀}:
		      \begin{itemize}
			      \item 有\*取值
			      \item 无\*地址
			      \item \&也是地址
		      \end{itemize}
	\end{itemize}
\end{bluebox}

\vspace{1cm}

% ======== 题目12: 二级指针 ========
\section*{题目12: 二级指针}

\textbf{原题目:}
设有定义:\texttt{int n=0,*p=\&n,**q=\&p;},则下列选项中正确的赋值语句是\_\_\_\_\_\_。
\begin{enumerate}
	\item[A)] \texttt{p=1;}
	\item[B)] \texttt{*q=2;}
	\item[C)] \texttt{q=p;}
	\item[D)] \texttt{*p=5;}
\end{enumerate}

\begin{bluebox}{答案与深度解析}
	\textbf{答案: D) \texttt{*p=5;}}

	% ======== 阶段一:核心认知框架 ========
	\textbf{【核心认知框架】}
	\begin{itemize}[leftmargin=*, nosep]
		\item \textbf{题目本质}: 考查二级指针的定义和赋值
		\item \textbf{关键区分点}: 指针的\textbf{值}与\textbf{指向}的区别
		\item \textbf{命题人意图}: 测试考生对多级指针的理解
	\end{itemize}

	% ======== 阶段二:二级指针解析 ========
	\textbf{【二级指针分析】}

	\textbf{① 定义解析}
	\begin{itemize}[leftmargin=*, nosep]
		\item \texttt{int n=0}: n是int变量,值为0
		\item \texttt{*p=\&n}: p是指针,指向n,\&n是n的地址
		\item \texttt{**q=\&p}: q是\textbf{二级指针},指向p,\&p是p的地址
	\end{itemize}

	\textbf{② 内存关系图}
	\begin{verbatim}
┌─────────┐         ┌─────────┐         ┌─────────┐
│    q    │────────→│    p    │────────→│    n    │
│ 0x1004  │         │  0x1008 │         │  0x100C │
│ 二级指针│         │  指针   │         │  int    │
└─────────┘         └─────────┘         └─────────┘
指向p的地址          指向n的地址          值为0
    \end{verbatim}

	\textbf{③ 访问层次}
	\begin{tabular}{|c|c|c|}
		\hline
		\textbf{表达式} & \textbf{类型}      & \textbf{含义}   \\
		\hline
		n            & int              & n的值(0)        \\
		\hline
		\&n          & int*             & n的地址          \\
		\hline
		p            & int*             & 指向n的指针(值=\&n) \\
		\hline
		*p           & int              & n的值(通过p访问)    \\
		\hline
		\&p          & int**            & p的地址          \\
		\hline
		q            & int**            & 指向p的指针(值=\&p) \\
		\hline
		*q           & int* | p的值(即\&n)                 \\
		\hline
		**q          & int | n的值(通过q访问)                 \\
		\hline
	\end{tabular}

	\textbf{④ 选项详细分析}

	\textbf{A) p=1}
	\begin{itemize}[leftmargin=*, nosep]
		\item \textbf{类型不匹配}: \texttt{p}是\texttt{int*}类型,1是\texttt{int}
		\item \textbf{含义}: 想把p指向地址1(危险操作)
		\item \textbf{编译警告}: 警告,但可编译(不安全)
	\end{itemize}

	\textbf{B) *q=2}
	\begin{itemize}[leftmargin=*, nosep]
		\item \texttt{*q}得到的是\texttt{p}的值(\texttt{\&n})
		\item 相当于: \texttt{\&n = 2}
		\item \textbf{错误}: 不能给地址\&n赋值2!
		\item \textit{注: *q是p,不是n}
	\end{itemize}

	\textbf{C) q=p}
	\begin{itemize}[leftmargin=*, nosep]
		\item \textbf{类型不匹配}: q是int**, p是int*
		\item \textbf{含义}: 试图把一级指针赋给二级指针
		\item \textbf{编译错误}: 类型不兼容
	\end{itemize}

	\textbf{D) *p=5}
	\begin{itemize}[leftmargin=*, nosep]
		\item \texttt{*p} = \texttt{n}的值的引用
		\item \textbf{含义}: n = 5
		\item \textbf{正确}: 给n赋值5
		\item \textbf{这是唯一合法的赋值}
	\end{itemize}

	% ======== 汇编理解 ========
	\textbf{⑤ 汇编级解释}
	\begin{lstlisting}[language=C]
// C: *p = 5;

// x86汇编:
mov eax, [ebp-12]  ; 加载p的值(= &n)
mov DWORD PTR [eax], 5  ; 向该地址写入5

// C: **q = 3;

// 汇编:
mov eax, [ebp-8]   ; 加载q的值(= &p)
mov eax, [eax]     ; 加载*q的值(= p = &n)
mov DWORD PTR [eax], 3  ; 写入3到n
    \end{lstlisting}

	% ======== 阶段三:应背诵知识点 ========
	\textbf{【应背诵知识点】}

	\begin{enumerate}[leftmargin=*, nosep]
		\item \textbf{多级指针定义}:
		      \begin{itemize}
			      \item \texttt{int *p}: 指向int的指针(一级指针)
			      \item \texttt{int **q}: 指向int*的指针(二级指针)
			      \item \texttt{int ***r}: 三级指针
		      \end{itemize}

		\item \textbf{访问规则}:
		      \begin{itemize}
			      \item \(k\)级指针需要\(k\)个\(*\)才能访问到值
			      \item \texttt{**q} = \texttt{*(*q)} = n
		      \end{itemize}

		\item \textbf{赋值规则}:
		      \begin{itemize}
			      \item 指针的值只能赋相同类型或兼容的指针
			      \item 不能把字面量赋给指针(除非NULL)
		      \end{itemize}
	\end{enumerate}

	% ======== 阶段四:实战策略 ========
	\textbf{【实战策略】}
	\begin{itemize}[leftmargin=*, nosep]
		\item \textbf{分析步骤}:
		      \begin{enumerate}
			      \item 画出指针链
			      \item 确定每个表达式的类型
			      \item 检查赋值兼容性
		      \end{enumerate}

		\item \textbf{常见错误}:
		      \begin{itemize}
			      \item 混淆*q和**q
			      \item 类型不匹配
			      \item 给地址赋值
		      \end{itemize}
	\end{itemize}
\end{bluebox}

\end{document}
